\chapter{Anthroponymie}\label{ch:noms}

Le relevé des noms n’est pas une prosopographie. Une forme peut appartenir à un défunt, à un commémorant, à un consul, à un vendeur d’emplacement ou à une personne citée dans une généalogie ; les changements de cas, les abréviations et les restitutions compliquent la reconnaissance. L’index des noms (chapitre~\ref{ch:index}) est produit automatiquement à partir du texte transmis : il retient les mots à initiale capitale, écarte les noms consulaires de l’année du texte et de l’année précédente (sauf lorsqu’ils désignent une autre personne), les noms de mois et de jours, et signale d’un astérisque les formes dont une lettre au moins est restituée. Il donne les formes telles qu’elles sont gravées, sans lemmatisation. Il reste un instrument de repérage.

\section{Gentilices et noms uniques}

La plupart des personnes commémorées portent un nom unique. Les gentilices sont rares et concentrés dans quelques familles ou quelques milieux. Hors des formules consulaires, \emph{Flavius} et \emph{Flavia} sont les plus fréquents (\nnums{19}, 20, 57, 59, 72, 149, 157, 217, 221 et 440), ce qui reflète la diffusion de ce nom, devenu presque un titre, dans le service impérial du \textsc{iv}\textsuperscript{e} siècle. \emph{Aurelius} et \emph{Aurelia} viennent ensuite (\nnums{6}, 19, 23, 111 et 229). Les autres gentilices n’apparaissent qu’une ou deux fois, souvent dans les textes de rang sénatorial : Cassia (\nnum{15}), Clodius (\nnum{36}), Roscia (\nnum{179}), Egnatia (\nnum{330}), Lucceia (\nnum{306}), Statilia (\nnum{42}), Cornificia (\nnum{77}), Aufidia (\nnum{449}), Valerius (\nnum{262}).

Ces effectifs ne mesurent ni la citoyenneté juridique ni la structure sociale de la population chrétienne : l’abandon du gentilice dans les épitaphes est un fait de pratique, non un indice de statut servile.

\section{Noms à signification chrétienne}

Un petit nombre de noms a une signification chrétienne directe ou probable : Adeodata (\nnum{297}) et peut-être Adeoda[tus] (\nnum{378}, restitué), Quodvultdeus (\nnum{388} ; \emph{Quodvuldeus} au \nnum{112}, en partie restitué), Theodora (\nnum{80}), Theodotus (\nnum{5}), Theophila (\nnum{110}), Anastasius et Anastasia (\nnums{43}, 195 et 197), Benedictus (\nnums{268} et 441), Martyrius (\nnum{139}), Pascasius (\nnums{235} et 400), Quiriace, c’est-à-dire Κυριακή (\nnum{298}). D’autres renvoient à des figures bibliques ou apostoliques (Petrus, \nnum{197} ; Paulus, \nnum{398} ; le Ioannes du \nnum{261}, gravé \emph{Ioannti}), ou au sabbat (Sabbatius, \nnums{22} et 386 ; Sambatia, \nnum{339}). Le nom \emph{Stercoria} (\nnum{108}, en partie restitué) appartient au groupe des noms d’humilité, dont Kajanto a montré la fréquence dans les inscriptions chrétiennes.\footnote{I.~Kajanto, \emph{Onomastic Studies in the Early Christian Inscriptions of Rome and Carthage}, Helsinki, 1963 (Acta Instituti Romani Finlandiae, II, 1) ; Ch.~Pietri, « Remarques sur l’onomastique chrétienne de Rome », dans \emph{L’onomastique latine}, Paris, 1977. Notices bibliographiques reprises de la bibliographie du volume ; passages non consultés.}

Ces noms ne sont pas des professions de foi. Leur analyse étymologique peut éclairer un horizon culturel ; mais le nom peut avoir été donné bien avant l’événement commémoré, relever d’une transmission familiale ou d’un usage devenu banal. La date d’une épitaphe atteste l’emploi d’un nom à cette date, non la date de son attribution, et l’usage d’un nom à étymologie chrétienne ne suffit pas à établir l’identité religieuse de son porteur. Inversement, la très grande majorité des défunts du corpus porte des noms latins ou grecs traditionnels (Felix, Victor, Leo, Laurentius, Ursus, Ianuaria, Crescentia…), sans signification chrétienne particulière.

\section{Doubles noms et désambiguïsation}

La formule \emph{Muscula quae et Galatea} (\nnum{151}) signale explicitement un double nom ; de même, \emph{Egnatia Susanna} est dite aussi \emph{Abita} (\nnum{330}), et Flavius Aurelius Casarius porte le surnom de \emph{Santulus} (\nnum{19}). Ces formulations doivent être distinguées des rapprochements entre deux fiches portant des noms voisins. L’identité d’une personne ne se déduit ni d’un nom commun ni d’une date proche : elle suppose la convergence d’indices documentaires. Les identifications proposées dans les dossiers commentés (Clodius Adelfius, Viventius, Proiecta) demeurent des hypothèses tant qu’elles n’ont pas été vérifiées dans la \emph{Prosopography of the Later Roman Empire} et la \emph{Prosopographie chrétienne du Bas-Empire} ; une référence à ces instruments désigne alors un contrôle à effectuer, non la preuve que ce contrôle a eu lieu.
